\section{Limitations and Future Work}\label{sec:limits}

\paragraph{Scope of models and SAEs.} The study covers two models and a finite set of SAE families: three in-house
architectures at one operating point (width 3{,}072, L0 = 128) and one data-scaled recipe in GPT-2, one public GPT-2
family, and one public SAE per layer at three layers of Gemma-2-2B. Our JumpReLU uses a sigmoid surrogate rather than
the published estimator and has many dead latents at layer 0, and a modern JumpReLU could not be trained to the target
sparsity in the scaling experiment. The conclusions describe these SAEs; they may not hold for other architectures,
widths or training budgets.

\paragraph{No controlled cross-model comparison.} GPT-2 and Gemma differ in model, SAE family, width, sparsity, hook
convention (\texttt{resid\_pre} versus \texttt{resid\_post}), data and context length. Gemma was evaluated at the
reconstruction and behavioural levels only. The study therefore supports a within-model replication of the direction
of one behavioural effect, not a comparison of the two models or of their SAEs.

\paragraph{Sparsity and capacity matching.} Architectures are matched on OWT sparsity, not on the sparsity they show
on task prompts, so architecture differences in information preservation remain confounded with task-level L0. The PCA
comparison depends on whether reconstruction error or the per-token active budget is equalised, and $k=128$ equalises
active dimensions, not parameters; only one of the three cells that are better at $k=128$ survives an approximate
post-hoc Holm correction. A TopK experiment trained at task-matched sparsity (E1) has been drafted but not approved or
run.

\paragraph{Reader dependence and incomplete causal evidence.} Information conclusions depend on the reader, and the
pre-registered truth criterion of the fine-tuning experiments used a frozen nonlinear probe that re-trained readers
disagree with. The causal evidence covers one persona-instruction effect of 0.02--0.04 nats at three sites of one model.
The patching design measures what the reconstruction preserves when the error term is kept, which is not mediation
through SAE latents. A mediated SAE-in-the-loop experiment has been pre-registered but has not been run, and the
fine-tuning comparisons at the layer-8 site can only resolve differences of about 0.05--0.1 in preservation share.

\paragraph{Validation-only results.} Only the OWT reconstruction table and the improvement-phase comparisons on OWT
and GoT are held-out test results. Every persona, causal, PCA-CM and Gemma result is a validation finding. A held-out
RepE-LM test protocol has been frozen, but it had not produced results when this version was written, so the layer-8
causal gain from scaling in particular awaits confirmation. Validation data were also used to compare methods, so
validation effect sizes may be optimistic. Three training seeds share the same data, so seed spread understates
uncertainty, and the fine-tuning criteria are applied per cell without multiplicity correction.

\paragraph{Ongoing work: R3.} R3 is a pre-registered robustness experiment that will test how Gemma Scope SAE width and
sparsity affect reconstruction and behavioural splicing at layers 5, 12 and 19, using the M6 sample and controls, and
how much the M6 values depend on the GPU's bf16 arithmetic. Its protocol and code are under review in a separate pull
request, and it has not produced a finished run. R3 is therefore
\evlabel{pre-registered only} and is not part of the completed evidence of this paper; none of our conclusions depends
on it.

\paragraph{Future work.} Beyond R3, the most direct extensions are to evaluate the frozen RepE-LM test protocol so that
the persona and causal results can be confirmed, to repeat the capacity-matched PCA comparison on the test split and on
the persona task, to train at task-matched sparsity, to run a mediated SAE-in-the-loop intervention, and to extend the
information and causal levels to Gemma-2-2B. Supervised low-rank subspaces, which carry more of the persona effect than
the full SAE reconstruction at layer 8, are a useful standard reference for future causal SAE evaluations.
